\documentclass[a4paper,11pt]{article}
\usepackage[margin=27mm]{geometry}
\usepackage{fontspec,amsmath,amssymb,amsthm,unicode-math}
\setmainfont{STIXGeneral}\setmathfont{STIX Two Math}
\newtheorem{theorem}{Theorem}\newtheorem{lemma}[theorem]{Lemma}\newtheorem{corollary}[theorem]{Corollary}
\newcommand{\Stab}{\operatorname{Stab}}\newcommand{\Aut}{\operatorname{Aut}}
\title{Intersections and faces in graph covers}
\author{Human Compression project}\date{Author edition; pending independent review}
\begin{document}\maketitle
\begin{abstract}
The fiber argument for simultaneous group actions also applies to an arbitrary
equivariant projection. For a rank-four cover whose endpoint faces stay connected
and whose middle faces lift isomorphically, this isolates the entire obstruction
to polytopality in the lifted vertex--facet sections. We prove the face and
thinness assertions explicitly. The application is the family of Leemans and
Toledo; the covering-space and group-action facts are standard, and no historical
priority is claimed for their reformulation here.
\end{abstract}

\section{An equivariant projection}
Let a group $W$ act on sets $Q$ and $X$, and let $\pi:Q\to X$ commute
with the action. Fix $q\in Q$, and put $x=\pi(q)$. For $S\leq W$, write
$D_S(q)=Sq\cap\pi^{-1}(x)$. Thus
\[
D_S(q)=(S\cap\Stab_W(x))q.
\]
\begin{lemma}\label{projection}
Suppose $C\leq A\cap B$ and $Ax\cap Bx=Cx$. Then
\[
Aq\cap Bq=Cq\quad\Longleftrightarrow\quad
D_A(q)\cap D_B(q)=D_C(q).
\]
\end{lemma}
\begin{proof}
The forward implication restricts the asserted equality to the fiber of $x$.
Conversely take $q'\in Aq\cap Bq$. Its image is $cx$ for some $c\in C$.
Then $c^{-1}q'$ lies in both orbits and in the fiber of $x$. The fiber
hypothesis puts it in $Cq$, and hence $q'\in Cq$. The opposite inclusion
is automatic. This is exactly the first-coordinate repair argument in
\emph{Intersections in simultaneous group actions}, Theorem~1: alternatively
apply that theorem to the invariant graph $\{(\pi(q),q):q\in Q\}$.
\end{proof}

For completeness, let $\Gamma$ be a connected graph with a voltage
$z(u,v)$ in a group $K$, satisfying $z(v,u)=z(u,v)^{-1}$. The cover has
vertices $(u,k)$ and edges from $(u,k)$ to $(v,z(u,v)k)$. Products along
walks are in traversal order on the left: the last edge acts last and
therefore appears on the left of the product. If $Y$ is a connected
subgraph and $u\in Y$, the voltages of its closed walks at $u$ form a
subgroup $H_Y\leq K$. Path lifting shows that the component of $(u,k)$
in the preimage of $Y$ meets the fiber of $u$ in $H_Yk$. Every component
meets that fiber, so its components are indexed by these cosets.
In particular its preimage is connected exactly when $H_Y=K$, and each
component maps isomorphically to $Y$ exactly when $H_Y=1$. These are the
usual elementary voltage-cover facts, proved here by following paths;
no finite enumeration is involved. If the voltage is trivial on a spanning
tree, every non-tree dart completes a fundamental closed walk whose voltage
is that dart's voltage. These fundamental walks generate all closed-walk
voltages, giving the source's reduced-voltage connectivity criterion.
For a cover of a maniplex, a distant-color square closes upstairs exactly
when its four dart voltages multiply to the identity. Consequently the
cover is a maniplex exactly when it is connected and every such square
has trivial voltage. This also recovers source Lemmas~3.3 and~3.4.

\section{A rank-four face criterion}
In the geometric sections all graphs are finite, as in the source.
A maniplex is a connected properly colored simple graph with one edge of every
color at each vertex, and commuting nonconsecutive color adjacencies.
Its flags are its vertices. An $i$-face is a component after deleting
color $i$; incidence means nonempty intersection. A cover preserves
colors and their unique adjacencies. Its restriction to any component
of a prescribed set of colors is a cover of the corresponding component
of the base: every color path lifts uniquely. All intersections below
retain the flags and edges belonging to every intersected face. In a
faithful maniplex this is also the induced graph on the common flags.

Faithfulness means that a flag is determined by all its faces. Thinness
means that every interval between incident faces whose ranks differ by two
has precisely two intermediate faces; improper endpoint faces are included.
We recall the elementary chain realization behind the face-poset convention.
Suppose a flag $x$ is already in prescribed faces of ranks at least $k$, and
we want also to meet an incident face $F_j$ of rank $j<k$. Choose
$y\in F_j\cap F_k$ and a path from $x$ to $y$ in $F_k$. Its colors avoid
$k$. Every color below $k$ commutes with every color above $k$, so reorder
the path into its lower-color part followed by its upper-color part. The
intermediate flag still belongs to all the prescribed higher faces; it
also belongs to $F_j$, since the final upper-color part cannot change that
face. Descending induction realizes every chain by a flag. The same argument
for three ranks proves transitivity of the face order.

If the maniplex is faithful and thin, this realization is bijective for
maximal chains and carries $i$-adjacency to changing just the rank-$i$ face.
Indeed $i$-adjacent flags lie in all the same other faces, and their
$i$-faces differ by faithfulness; thinness says these are the only two
completions. Thus a faithful thin maniplex is polytopal exactly when all
face-chain intersections are connected. This is the usual connected-
intersection characterization of Garza-Vargas and Hubard; the preceding
argument gives precisely the direction and hypotheses used here.

\begin{lemma}\label{smallrank}
Every faithful thin rank-three maniplex is polytopal. In a faithful
rank-four maniplex each middle face is a prism: a single edge times an
even alternating polygon. If the rank-four maniplex is also thin, every
chain containing a middle face has connected intersection.
\end{lemma}
\begin{proof}
In any faithful maniplex, $i$-adjacent flags cannot belong to the same
$i$-face: their other faces already agree, and agreement of the remaining
face would contradict faithfulness. For a faithful thin rank-three
maniplex, an intersection of two faces consists of the two flags
completing that chain, joined by the remaining color. Three faces have
one common flag, a singleton face is connected by definition, and the
whole graph is connected. The characterization just proved applies.
This recovers the full scope of Leemans--Toledo Lemma~2.6.

For a rank-four $1$-face, consider one colors-$2,3$ polygon. Color~0
commutes with both of these colors, so it maps this polygon isomorphically
to another. The polygons are distinct: the first is contained in one
$0$-face, and a $0$-neighbor cannot remain in that face. Connectivity
then says these two polygons exhaust the $1$-face, paired by color~0.
This proves the prism assertion; the $2$-face is dual. In particular a
$1$-face meets exactly two $0$-faces, and a $2$-face meets exactly two
$3$-faces. No thinness assumption was used here.

Assume now thinness. A chain of at least three proper faces has either
one flag or the two flags of the missing-color edge, as above. For a
pair of faces it remains to examine five pairs of ranks. An incident
$0,1$ pair intersects in exactly one of the polygon layers just described;
the $2,3$ pair is dual. For a $1,2$ pair there are at most two choices
of vertex and two of facet, hence at most four flags by faithfulness.
The colors-$0,3$ square through any common flag supplies four, so is the
whole intersection. For a $0,2$ pair thinness gives exactly two possible
intermediate $1$-faces, and the prism description gives two possible
facets. Again there are at most four flags, and the colors-$1,3$ square
supplies all four. The $1,3$ case is dual. These are all pairs containing
a middle face. Singleton faces and the larger chains already considered
finish the proof.
\end{proof}

\begin{theorem}\label{faces}
Let $P$ be a faithful thin rank-four maniplex, and let $\pi:Q\to P$ be a
connected maniplex cover. Assume:
\begin{enumerate}
\item the full preimage of every $0$-face and every $3$-face is connected;
\item every component over a $1$-face or a $2$-face maps isomorphically
onto that face.
\end{enumerate}
Then $Q$ is faithful and thin. Every chain containing a middle face has
a connected nonempty intersection, mapped isomorphically to the
intersection of its image faces. Consequently $Q$ is polytopal if and
only if the preimage of every incident vertex--facet intersection in
$P$ is connected.
\end{theorem}
\begin{proof}
Endpoint faces of $Q$ are the full preimages of endpoint faces of $P$.
Fix middle faces $H_1,H_2$ of $Q$ that meet, and write $F_i=\pi(H_i)$.
The base intersection $F_1\cap F_2$ is connected. Starting at a flag
of $H_1\cap H_2$, every path in that base intersection lifts to a path
in both faces, because it uses neither color~1 nor color~2. Since each
$H_i$ maps isomorphically, there is exactly one lift of each base flag
in either $H_i$. Their restrictions therefore agree over the entire
intersection. Thus $H_1\cap H_2$ maps isomorphically onto $F_1\cap F_2$.
This establishes the projection-of-intersection equality in this
specific situation; such equality is not assumed for arbitrary covers.

Now consider any chain containing at least one middle face. Its endpoint
faces impose exactly the conditions of belonging to their full preimages.
If both middle faces occur, the preceding paragraph identifies their
intersection; if only one occurs, its own projection is an isomorphism.
In either case restricting that isomorphism to the endpoint conditions
identifies the whole intersection with its base-chain intersection.
The latter is connected by Lemma~\ref{smallrank} and nonempty by chain realization. For a
maximal chain it is a singleton. In particular each flag is determined
by its incident faces, which is faithfulness.

We verify thinness instead of inferring it just from injectivity of a
projection. Fix a chain missing one rank $i$. It still contains at least
one middle face. Its intersection is, by the preceding argument,
isomorphic to the base intersection for the image chain. That base
intersection consists of precisely the two flags of an $i$-edge.
Both lifted flags are in all the prescribed faces. Their respective
$i$-faces complete the chain, and these faces are distinct, since
otherwise faithfulness would identify the two flags. Conversely, every
completion has a flag in its full intersection by the maximal-chain
argument, so its $i$-face is one of these two. Every codimension-one chain
therefore has precisely two completions. Extending any incident pair
of ranks $i-1,i+1$ to such a chain shows the diamond condition: any
intermediate $i$-face is comparable with every chosen lower or upper
face by transitivity of the face order, and hence is one of those two.

The only remaining chains have no middle face. A singleton face is
connected; a chain consisting of a vertex and a facet has intersection
exactly the full preimage of their base intersection. The empty chain
has intersection $Q$, which is connected. The connected-intersection characterization with its faithful and thin
hypotheses proves the last assertion.
\end{proof}

\begin{corollary}\label{double}
Let $P$ be a faithful rank-four maniplex whose $0$- and $3$-faces are
nonbipartite. Its canonical double cover, with adjacency
$(u,\epsilon)^i=(u^i,\epsilon+1)$, is connected and faithful and fails
the connected-intersection property. It has the same numbers of vertices
and facets and twice the numbers of edges and two-faces. If $P$ is thin,
so is its double cover, which is then nonpolytopal. If $P$ is regular,
so is its double cover, and its automorphism group is $\Aut(P)\times C_2$.
\end{corollary}
\begin{proof}
Along a lifted walk the second coordinate records its length modulo2.
The double cover of a connected graph is therefore connected precisely
when that graph has an odd closed walk; otherwise it consists of two
isomorphic copies. Endpoint faces have connected preimages by hypothesis.
The middle faces are bipartite prisms by Lemma~\ref{smallrank}, so each
has two isomorphic lifts. This proves connectedness and the face counts.
If two flags upstairs lie in all the same faces, their projections do so
in $P$ and hence coincide. Both upstairs flags are in one lifted middle
face, where projection is injective, so they coincide as well. This proves
faithfulness without a thinness assumption. When $P$ is thin,
Theorem~\ref{faces} proves thinness of the cover.

Choose incident endpoint faces in $P$. Their intersection has only
colors~1,2, with one edge of each color at every flag. It is a disjoint
union of even alternating polygons. If there is more than one component,
its preimage is disconnected; if there is one, its canonical double
already has two components. Since full endpoint preimages are faces,
this disconnected preimage is an actual face intersection upstairs.
Thus the double cover fails the connected-intersection property. If it
is thin, its proved faithfulness makes it nonpolytopal.

Each automorphism $\alpha$ of $P$ acts by
$(u,\epsilon)\mapsto(\alpha(u),\epsilon)$, commuting with the deck flip.
If $P$ is regular, this gives a transitive subgroup $\Aut(P)\times C_2$.
A color-preserving automorphism of a connected maniplex is determined by
one flag, so the transitive semiregular subgroup is the entire automorphism
group. This proves every clause of source Lemma~3.8 on its stated faithful
base, despite the more general thin-base Lemma~3.5 being false.
\end{proof}

\section{Lifting an automorphism}
The standard sufficient lift condition has a useful version with a change
of coordinates in each fiber. Suppose voltages take values in an additive
abelian group $K$, $\phi$ is a base automorphism, $f\in\Aut(K)$, and
$h$ is a function on the base vertices. If every dart $u v$ satisfies
\[
z(\phi u,\phi v)=f(z(u,v))+h(v)-h(u),
\]
then $(u,k)\mapsto(\phi u,f(k)+h(u))$ is a lift of $\phi$.
It is bijective on each fiber; the displayed equation says exactly that
the images of adjacent vertices are adjacent. This proves the claim.
For arbitrary multiplicative $K$ and $h=0$, the same direct calculation
proves $(u,k)\mapsto(\phi u,f(k))$ whenever
$z(\phi u,\phi v)=f(z(u,v))$. This is source Lemma~5.5, attributed there
to Malni\v c--Nedela--\v Skoviera and Poto\v cnik--Toledo. We use only
this elementary sufficient condition, not an unstated converse.

\section{The necessary restriction on the base}
The faithful hypothesis in Theorem~\ref{faces} cannot simply be dropped.
Leemans--Toledo Lemma~3.5 assumes only a thin base with the same two
conditions on face preimages. The application gives a complete explicit
32-to-64-flag counterexample. Its base is thin but unfaithful; its cover
is faithful but not thin. Thus the exact general statement is disposed
of by a counterexample, while Lemma~\ref{smallrank} and
Theorem~\ref{faces} supply a sufficient repair covering the source's
faithful construction. No claim about the source's principal theorem
is inferred from failure of that auxiliary lemma.
\end{document}
